\section{Every order derivative on the odd Lerch lattice}
\label{lj:sec:main}

The report \emph{Continuous Depth Lerch} asks in its final question Q8
for the spectral derivatives of its continuous-depth family at odd
positive integers \cite{Continuous}. Its formula \texttt{cdl:eq:oddF} gives the values there,
but does not permit differentiation in the discrete index. We resolve
that question in the precise form requested: an explicit finite formula
for every order derivative and every mixed depth derivative, including
the reflected spectral coordinates. We also give normally convergent
series for all those coordinates. Thus the answer includes a constructive
evaluation procedure at the negative reflected orders, where the original
alternating series need not converge.

The analytic inputs are the source report's proved centered reflection
formula \texttt{cdl:eq:completed}, Euler's classical transformation of
series, and the usual Gamma and Hurwitz-zeta identities \cite{DLMFEuler,DLMFGamma,DLMFZeta}. The new claim is
the following continuation of the inspected source question, not priority
for those classical inputs or for the special cases recovered below.
Euler's transformation is recorded in DLMF \S3.9(ii); the normally
convergent, parameter-differentiated use needed here is proved explicitly.

\subsection{Functions, branches, and normalization}

For $\re\rho>0$, put
\begin{align}
 f_\rho(t)&=(2\cosh(t/2))^{-\rho},\qquad t\geq0,
 \label{lj:eq:f}\\
 Z_\rho(s)&=\frac1{\Gamma(s)}
       \int_0^\infty t^{s-1}f_\rho(t)\,\dd t,
       \qquad \re s>0.
 \label{lj:eq:Z}
\end{align}
The logarithm in $f_\rho$ is real on the positive real axis. The function
$Z_\rho(s)$ is entire in $s$: subtracting finitely many Taylor terms of
$f_\rho$ near zero gives its meromorphic Mellin transform, and the zeros
of $1/\Gamma(s)$ cancel its poles. The same subtractions on compact
parameter sets establish joint holomorphy for $\re\rho>0$. In the
absolutely convergent region $\re s>\re\rho$,
\begin{equation}
 Z_\rho(s)=\sum_{k=0}^\infty
     \frac{(-1)^k(\rho)_k}{k!(k+\rho/2)^s}.
 \label{lj:eq:Zseries}
\end{equation}
Powers of $k+\rho/2$ use the logarithm on the right half-plane.

Define
\begin{align}
 b_\rho(x)&=\frac{(2\pi)^{\rho-2}}{\Gamma(\rho)}
       \Gamma\left(\frac\rho2+\frac{ix}{2\pi}\right)
       \Gamma\left(\frac\rho2-\frac{ix}{2\pi}\right),
 \label{lj:eq:b}\\
 F_\rho(s)&=\frac1{\Gamma(s)}
       \int_0^\infty x^{s-1}b_\rho(x)\,\dd x,
       \qquad \re s>0.
 \label{lj:eq:F}
\end{align}
This is $\mathcal F_{\rho,1/2}(s;0)$ in the continuous-depth report.
Its entire continuation follows from the same Mellin-subtraction argument.
For complex $\rho$, the Gamma product is a product of two functions,
not a squared modulus.

The centered reflection identity is
\begin{equation}
 F_\rho(s)=\frac{\pi(2\pi)^{\rho+s-2}}
                 {\Gamma(s)\sin(\pi s/2)}Z_\rho(1-s).
 \label{lj:eq:reflection}
\end{equation}
Apparent singularities are removed in the combined expression. For
completeness, the beta integral gives the classical Fourier formula \cite{RamanujanFourier}
$\widehat f_\rho(v)=
\Gamma(\rho/2+iv)\Gamma(\rho/2-iv)/\Gamma(\rho)$, where the Fourier
transform is $\int_{\mathbb R}e^{ivt}f_\rho(t)\,\dd t$. On
$0<\re s<1$, insert a decaying exponential in the cosine Mellin
integral, use Fubini, and then dominated convergence to obtain
\[
 \int_0^\infty v^{s-1}\widehat f_\rho(v)\,\dd v
 =2\Gamma(s)\cos(\pi s/2)
       \int_0^\infty t^{-s}f_\rho(t)\,\dd t.
\]
The regulator can be removed on the left because the Gamma product
decays exponentially. Rescaling gives
$F_\rho(s)=2(2\pi)^{\rho+s-2}\Gamma(1-s)
\cos(\pi s/2)Z_\rho(1-s)$; Euler reflection for Gamma gives
\eqref{lj:eq:reflection}. Entire continuation gives the stated global
interpretation. This reproduces the source input rather than replacing
it with a formal exchange at a pole.

We use $\partial_s^q$ for order derivatives and $\partial_\rho^r$ for
depth derivatives. The complete exponential Bell polynomials are fixed by
\begin{equation}
 \exp\left(\sum_{j\geq1}d_j\frac{v^j}{j!}\right)
 =\sum_{q\geq0}\mathcal B_q(d_1,\ldots,d_q)\frac{v^q}{q!},
 \qquad \mathcal B_0=1.
 \label{lj:eq:Bell}
\end{equation}
Write $H_N^{(j)}=\sum_{k=1}^N k^{-j}$ and $H_0^{(j)}=0$.

\subsection{The finite formula for all mixed jets}

\begin{theorem}[Odd-lattice order and depth derivatives]
\label{lj:thm:oddjets}
Let $m,r,q\geq0$ be integers, put $N=2m$, $n=N+1$, and set
\begin{align}
 L&=\log(2\pi),\qquad
 C_{m,\rho}=\frac{(-1)^m\pi(2\pi)^{\rho+N-1}}{N!},
 \label{lj:eq:C}\\
 \kappa_1&=L+\gamma-H_N,\nonumber\\
 \kappa_j&=(j-1)!\left[
       \left(1-\frac{1+(-1)^j}{2^j}\right)\zeta(j)
       +(-1)^jH_N^{(j)}\right],\qquad j\geq2.
 \label{lj:eq:kappa}
\end{align}
For every $\re\rho>0$,
\begin{equation}
 \boxed{\displaystyle
 \partial_\rho^r\partial_s^qF_\rho(n)
 =C_{m,\rho}\sum_{j=0}^q(-1)^j\binom qj
       \mathcal B_{q-j}(\kappa_1,\ldots,\kappa_{q-j})
       (\partial_\rho+L)^r\partial_s^jZ_\rho(-N).}
 \label{lj:eq:alljets}
\end{equation}
The operator on the right is explicitly
$\sum_{h=0}^r\binom rh L^{r-h}\partial_\rho^h$.
No differentiation with respect to $m$ occurs.
\end{theorem}

\begin{proof}
The pole-free local identity from which every derivative follows is
\begin{equation}
 F_{\rho+u}(n+v)=C_{m,\rho}e^{Lu}
       R_m(v)Z_{\rho+u}(-N-v),
 \qquad
 R_m(v)=e^{Lv}\frac{\Gamma(n)}{\Gamma(n+v)}
                  \sec(\pi v/2).
 \label{lj:eq:local}
\end{equation}
It holds for $u$ sufficiently small and $|v|<1$, by substituting in
\eqref{lj:eq:reflection} and using
$\sin(\pi(n+v)/2)=(-1)^m\cos(\pi v/2)$.
All factors in \eqref{lj:eq:local} are holomorphic there.

The Gamma contribution to the $j$th logarithmic derivative of $R_m$
is $-\psi^{(j-1)}(n)$. For $j\geq2$,
\[
 \psi^{(j-1)}(n)=(-1)^j(j-1)!
                         \bigl(\zeta(j)-H_N^{(j)}\bigr).
\]
The product for cosine gives the convergent expansion
\[
 \log\sec(\pi v/2)
   =\sum_{k\geq1}\frac{(1-2^{-2k})\zeta(2k)}{k}v^{2k},
   \qquad |v|<1.
\]
Combining these expansions gives exactly \eqref{lj:eq:kappa} and
$R_m^{(a)}(0)=\mathcal B_a(\kappa_1,\ldots,\kappa_a)$.
Now apply the ordinary multivariate Leibniz rule to the holomorphic
identity \eqref{lj:eq:local}. The argument $-N-v$ supplies $(-1)^j$,
while $e^{Lu}$ supplies $(\partial_\rho+L)^r$. This proves
\eqref{lj:eq:alljets} with every combinatorial coefficient fixed.
\end{proof}

For example, putting $Z^{[j]}=\partial_s^jZ_\rho(-N)$, the first two
order derivatives are
\begin{align}
 F_\rho'(n)&=C_{m,\rho}(\kappa_1Z^{[0]}-Z^{[1]}),
 \label{lj:eq:first}\\
 F_\rho''(n)&=C_{m,\rho}
       \bigl((\kappa_1^2+\kappa_2)Z^{[0]}
                      -2\kappa_1Z^{[1]}+Z^{[2]}\bigr),
 \label{lj:eq:second}\\
 \kappa_2&=\tfrac12\zeta(2)+H_N^{(2)},\qquad
 \kappa_3=2\bigl(\zeta(3)-H_N^{(3)}\bigr).\nonumber
\end{align}

To recover the already known value, define
$U_m(\rho)=[z^{2m}]\sec(z)^\rho$. The even Taylor expansion of
$f_\rho$ gives
\begin{equation}
 Z_\rho(-2m)=(-1)^m(2m)!2^{-\rho-2m}U_m(\rho),
 \qquad
 F_\rho(2m+1)=\frac{\pi^{\rho+2m}}2U_m(\rho).
 \label{lj:eq:valuecheck}
\end{equation}
These are the $q=0$ inputs. The terms $Z^{[j]}$ with $j>0$ in
\eqref{lj:eq:alljets} are additional data; discarding them would give
an incorrect spectral derivative.

\subsection{A convergent elementary series for every reflected jet}

The next theorem makes the reflected coordinates in
\eqref{lj:eq:alljets} independently computable. In particular, it avoids
using the reflection equation itself as a numerical definition of the
coordinates that are supposed to verify that equation.

\begin{theorem}[Normally convergent Euler series at every order]
\label{lj:thm:Euler}
For $\re\rho>0$ and $s\in\mathbb C$, define
\begin{equation}
 A_k(\rho,s)=\frac{(\rho)_k}{k!(k+\rho/2)^s},\qquad
 D_\ell(\rho,s)=\sum_{k=0}^\ell(-1)^k\binom\ell k A_k(\rho,s).
 \label{lj:eq:D}
\end{equation}
Then
\begin{equation}
 \boxed{\displaystyle
 Z_\rho(s)=\sum_{\ell=0}^\infty\frac{D_\ell(\rho,s)}{2^{\ell+1}}.}
 \label{lj:eq:Euler}
\end{equation}
This series is absolutely and locally normally convergent throughout
$\{\re\rho>0\}\times\mathbb C$. Every fixed mixed derivative in
$\rho,s$ is obtained by differentiating its finite inner sums termwise.
\end{theorem}

\begin{proof}
For $z\in\Omega=\mathbb C\setminus[1,\infty)$, initially $\re s>0$,
let
\begin{equation}
 A(\rho,s;z)=\frac1{\Gamma(s)}\int_0^\infty
        t^{s-1}e^{-\rho t/2}(1-ze^{-t})^{-\rho}\,\dd t.
 \label{lj:eq:Aint}
\end{equation}
Use the principal logarithm of $1-ze^{-t}$. It is available on
$\Omega$, since $1-ze^{-t}$ cannot lie on the nonpositive real axis
there. On compact subsets, the integrand is bounded at zero after the
factor $t^{s-1}$ is removed and decays exponentially at infinity.
Thus \eqref{lj:eq:Aint} is jointly holomorphic in this initial domain.
Its power series at $z=0$ is $\sum_{k\geq0}A_k(\rho,s)z^k$.

The coefficient identity
\begin{equation}
 \left(z\frac\partial{\partial z}+\frac\rho2\right)
 A(\rho,s+1;z)=A(\rho,s;z)
 \label{lj:eq:Aladder}
\end{equation}
holds first in $|z|<1$ and then throughout $\Omega$. Iterating it from
$\re(s+M)>0$ continues $A$ to every $s\in\mathbb C$. The definitions
agree on overlaps by the same identity, so this is a joint holomorphic
continuation. At $z=-1$, it equals $Z_\rho(s)$.

Set
\[
 B(\rho,s;v)=\frac1{1-v}
              A\left(\rho,s;\frac{-v}{1-v}\right).
\]
If $|v|<1$, then $-v/(1-v)$ is outside $[1,\infty)$: a point of that
cut would require $v\in(1,\infty)$. Hence $B$ is holomorphic for
$|v|<1$. Its Taylor coefficient is $D_\ell(\rho,s)$, by expanding
$(1-v)^{-k-1}$ in the defining power series near zero.
Taking $v=1/2$ yields $B(\rho,s;1/2)=2Z_\rho(s)$ and proves
\eqref{lj:eq:Euler}.

For a compact parameter set, choose any fixed radius $r$ with
$1/2<r<1$. Cauchy's coefficient estimate on $|v|=r$ gives
$|D_\ell|\leq C r^{-\ell}$ uniformly on that set. It proves absolute
normal convergence of \eqref{lj:eq:Euler}. Apply the same argument on
a slightly enlarged compact parameter neighborhood to justify every
fixed mixed derivative. This argument proves convergence even when
\eqref{lj:eq:Zseries} is a divergent ordinary series.
\end{proof}

All finite derivative coefficients admit a harmonic expression. Put
$a_k=k+\rho/2$ and
\[
 h_k^{(j)}(\rho)=\sum_{\ell=0}^{k-1}(\rho+\ell)^{-j},\qquad
 d_j=(-1)^{j-1}(j-1)!
       \left(h_k^{(j)}(\rho)-\frac{s}{2^ja_k^j}\right).
\]
Then $\partial_\rho^r A_k=A_k\mathcal B_r(d_1,\ldots,d_r)$.
For mixed derivatives a particularly safe finite generating formula is
\begin{equation}
 \frac{A_k(\rho+u,s+v)}{A_k(\rho,s)}
 =e^{-v\log a_k}
       \prod_{\ell=0}^{k-1}\left(1+\frac{u}{\rho+\ell}\right)
       \left(1+\frac{u}{2a_k}\right)^{-s-v}.
 \label{lj:eq:mixedcoeff}
\end{equation}
Only the requested bidegree is needed. This also distinguishes a mixed
derivative from the product of two independent logarithmic derivatives.

At the reflected integer $s=-N$, $N\geq0$, the formula is especially
explicit:
\begin{equation}
\begin{split}
 \partial_\rho^r\partial_s^q A_k(\rho,-N)
 &=(-1)^q r!\frac{(\rho)_k}{k!}a_k^N[u^r]
       \prod_{\ell=0}^{k-1}\left(1+\frac{u}{\rho+\ell}\right)
       \left(1+\frac{u}{2a_k}\right)^N\\
 &\hspace{13mm}\cdot
       \left(\log a_k+\log\left(1+\frac{u}{2a_k}\right)\right)^q.
\end{split}
 \label{lj:eq:negativecoeff}
\end{equation}
For positive rational $\rho$, every inner coefficient here is a rational
linear combination of powers of logarithms of positive rational numbers.
Equations \eqref{lj:eq:negativecoeff}, \eqref{lj:eq:Euler}, and
\eqref{lj:eq:alljets} therefore constitute a complete ordinary convergent
series evaluation of all the reflected mixed jets in the theorem.
The outer Euler series need not terminate. At a positive integer $\rho$
and $q=r=0$, its inner sequence is polynomial at negative integer $s$,
so finite differences do terminate; differentiation can remove that
termination.

\subsection{Independent Gamma-moment realization}

The positive-half-line Gamma integral gives a different exact
representation for checking the lattice formulas. Put
\begin{equation}
 a_j(\rho,x)=\delta_{j1}\log(2\pi)-\psi^{(j-1)}(\rho)
  +2^{-j}\left\{
   \psi^{(j-1)}\left(\frac\rho2+\frac{ix}{2\pi}\right)
  +\psi^{(j-1)}\left(\frac\rho2-\frac{ix}{2\pi}\right)\right\}.
 \label{lj:eq:Gammajet}
\end{equation}
For the order derivatives define $c_1(x;s)=\log x-\psi(s)$ and
$c_j(x;s)=-\psi^{(j-1)}(s)$ for $j\geq2$. Direct differentiation of
\eqref{lj:eq:F} gives, for $\re s>0$,
\begin{equation}
\begin{split}
 \partial_\rho^r\partial_s^qF_\rho(s)
  =\frac1{\Gamma(s)}\int_0^\infty
      x^{s-1}b_\rho(x)
      \mathcal B_r(a_1(\rho,x),\ldots,a_r(\rho,x))\qquad\qquad\\[-2pt]
      {}\cdot\mathcal B_q(c_1(x;s),\ldots,c_q(x;s))\,\dd x.
\end{split}
 \label{lj:eq:momentjets}
\end{equation}
This is an ordinary absolutely convergent integral. Near zero, only
fixed powers of $\log x$ occur. At infinity, the Gamma product decays
as $e^{-x/2}$ times a power of $x$; its parameter derivatives introduce
only logarithmic or polynomial factors on compact parameter sets.
These observations justify every differentiation. The left side of
\eqref{lj:eq:alljets} can therefore be evaluated through
\eqref{lj:eq:momentjets} while its right side is evaluated through
the independent Euler series.

\subsection{Stieltjes coordinates at depth two}

Let the generalized Stieltjes constants have the convention
\[
 \zeta(1+v,a)=\frac1v+
       \sum_{j\geq0}\frac{(-1)^j\gamma_j(a)}{j!}v^j,
       \qquad \re a>0.
\]
An upper parenthesized index on $\gamma_j$ below is an argument
derivative. Set $E_a(s)=s\zeta(s+1,a)$, with the removable value
$E_a(0)=1$.

\begin{proposition}[Stieltjes transport to every positive order]
\label{lj:prop:Stieltjes}
For integers $n\geq1,q\geq0$, put $N=n-1$ and
$\lambda_j=(-1)^j(j-1)!H_N^{(j)}$. Then
\begin{equation}
 \boxed{\displaystyle
 E_a^{[q]}(n)=\frac{(-1)^n}{N!}
   \sum_{j=0}^q(-1)^j\binom qj
       \mathcal B_{q-j}(\lambda_1,\ldots,\lambda_{q-j})
       \gamma_j^{(n)}(a).}
 \label{lj:eq:Stieltjes}
\end{equation}
In particular,
\begin{equation}
 E_a^{[q]}(1)=(-1)^{q+1}\gamma_q'(a).
 \label{lj:eq:Stieltjesone}
\end{equation}
At $a=1/2$, this gives $\partial_s^qF_2(n)$.
\end{proposition}

\begin{proof}
The shift derivative of Hurwitz zeta gives
$\partial_a^n\zeta(1+v,a)=(-1)^n(1+v)_n\zeta(n+1+v,a)$.
The pole $1/v$ is independent of $a$, so
\begin{equation}
 E_a(n+v)=\frac{(-1)^n}{(1+v)_N}
       \sum_{j\geq0}\frac{(-1)^j\gamma_j^{(n)}(a)}{j!}v^j.
 \label{lj:eq:Stieltjesgerm}
\end{equation}
The logarithmic Taylor coefficients of $N!/(1+v)_N$ are $\lambda_j$.
Coefficient extraction proves \eqref{lj:eq:Stieltjes}.
Finally, $b_2(x)=x/(2\sinh(x/2))$, so its Mellin transform gives
$F_2(s)=s\zeta(s+1,1/2)$ directly. This proves the last assertion.
\end{proof}

This proposition is an explicit form of the familiar parameter-derivative
tower. Its role here is to transport the new lattice-jet formula into
the repository's Stieltjes convention; it is not presented as a new
Hurwitz-zeta identity. For instance, with $N=2m$,
\[
 F_2'(2m+1)=\frac{
       \gamma_1^{(2m+1)}(1/2)
        +H_{2m}\gamma_0^{(2m+1)}(1/2)}{(2m)!}.
\]
Together with \eqref{lj:eq:first}, this is an exact relation between
negative reflected $Z_2$ order derivatives and positive-order Stieltjes
argument derivatives.

\subsection{Gamma-weighted polylogarithm order derivatives}

Put
\begin{align}
 G_\rho(x)&=
   \Gamma\left(\frac\rho2+\frac{ix}{2\pi}\right)
   \Gamma\left(\frac\rho2-\frac{ix}{2\pi}\right),\nonumber\\
 P_s(x)&=\frac{\operatorname{Li}_s(ie^{-x/2})
                    -\operatorname{Li}_s(-ie^{-x/2})}{i},\nonumber\\
 I_\rho(s)&=\int_{\mathbb R}G_\rho(x)P_s(x)\,\dd x.
 \label{lj:eq:I}
\end{align}
The polylogarithms use their principal branches. Their arguments stay
off the positive real cut. The following formula covers all independent
order derivatives, including orders for which the integrand is no longer
an elementary logarithm or a classical fixed-order polylogarithm.

\begin{theorem}[All polylogarithmic order jets under a Gamma weight]
\label{lj:thm:polylog}
For $\re\rho>0$, the integral $I_\rho(s)$ is entire in $s$, jointly
holomorphic in its parameters, and
\begin{equation}
 I_\rho(s)=\Gamma(\rho)(2\pi)^{2-\rho}2^{1-s}F_{\rho+1}(s)
 =\frac{4\Gamma(\rho)\pi^{s+2}}
        {\Gamma(s)\sin(\pi s/2)}Z_{\rho+1}(1-s),
 \label{lj:eq:Imaster}
\end{equation}
with all apparent poles removed in the combined expression.
For $N=2m$, $n=N+1$, replace only the first cumulant in
\eqref{lj:eq:kappa} by
$\widetilde\kappa_1=\log\pi+\gamma-H_N$, and leave
$\widetilde\kappa_j=\kappa_j$ for $j\geq2$. Then
\begin{equation}
 \partial_s^qI_\rho(n)=
 \frac{4(-1)^m\Gamma(\rho)\pi^{n+2}}{N!}
   \sum_{j=0}^q(-1)^j\binom qj
    \mathcal B_{q-j}(\widetilde\kappa_1,\ldots,
                    \widetilde\kappa_{q-j})
    \partial_s^jZ_{\rho+1}(-N).
 \label{lj:eq:Ijets}
\end{equation}
Every mixed depth derivative follows by the finite rule
\begin{equation}
 \partial_\rho^r\{\Gamma(\rho)Y(\rho+1)\}
 =\Gamma(\rho)\sum_{h=0}^r\binom rh
     \mathcal B_{r-h}(\psi(\rho),\psi'(\rho),\ldots)
     Y^{(h)}(\rho+1).
 \label{lj:eq:Igammajets}
\end{equation}
All these derivatives pass through the ordinary integral in
\eqref{lj:eq:I}.
\end{theorem}

\begin{proof}
The root filter in the continuous-depth report, valid first by a
convergent series and then by analytic continuation, is
\[
 \mathcal F_{1,1/2}(s;-x)=2^{s-1}e^{-x/2}P_s(x).
\]
Convolve this with
$K_\rho(x)=(2\pi)^{\rho-2}e^{x/2}G_\rho(x)/\Gamma(\rho)$.
The source's entire-order addition law gives $F_{\rho+1}(s)$ and
therefore the first equality in \eqref{lj:eq:Imaster}. The second is
\eqref{lj:eq:reflection}, with the powers of $2$ and $\pi$ collected.
The ordinary integral is absolutely convergent locally uniformly in
all the stated parameters. One explicit justification uses the vertical
contour for $\mathcal F_{1,1/2}$ on any line $1/2<c<1$, giving
$|P_s(x)|\leq C e^{(1/2-c)x}$ on compact order sets. The Gamma weight
then decays integrably at both ends. Fixed order derivatives add powers
of the logarithm of the contour variable, and depth derivatives of the
Gamma weight add only polynomial or logarithmic factors. The same
integrable majorants persist on slightly larger compact sets, proving
joint holomorphy and passage of all derivatives.

The local germ at $s=n+v$ now contains $e^{v\log\pi}$ in place of
$e^{v\log(2\pi)}$. The proof of Theorem~\ref{lj:thm:oddjets} gives
\eqref{lj:eq:Ijets}. Finally, the logarithmic derivatives of
$\Gamma(\rho)$ are $\psi^{(j-1)}(\rho)$, so Bell's formula and Leibniz
give \eqref{lj:eq:Igammajets}.
\end{proof}

On the integral side of \eqref{lj:eq:Igammajets}, the weight replacing
$G_\rho$ is
\[
 G_\rho(x)\mathcal B_r(g_1(\rho,x),\ldots,g_r(\rho,x)),\qquad
 g_j(\rho,x)=2^{-j}\left\{
 \psi^{(j-1)}\left(\frac\rho2+\frac{ix}{2\pi}\right)
 +\psi^{(j-1)}\left(\frac\rho2-\frac{ix}{2\pi}\right)\right\}.
\]
Thus the theorem is a family of explicit ordinary integrals containing
polygamma factors and spectral derivatives of polylogarithms. Theorem
\ref{lj:thm:Euler} evaluates their reflected coordinates by convergent
series, including the mixed derivatives appearing in
\eqref{lj:eq:Igammajets}.

At $\rho=1$ and $s=1$, $G_1(x)=\pi\operatorname{sech}(x/2)$.
The first derivative has the concrete forms
\begin{align}
 \int_{\mathbb R}\frac{\pi}{\cosh(x/2)}
       \left.\partial_sP_s(x)\right|_{s=1}\,\dd x
 &=2\pi\gamma_1'(1/2)-\pi^3\log2
 \label{lj:eq:exampleStieltjes}\\
 &=\pi^3\left(\gamma+\log\pi+
                  \frac43\log2+12\zeta'(-1)\right).
 \label{lj:eq:exampleZeta}
\end{align}
Indeed, the first equality follows from
$I_1(s)=2\pi2^{1-s}E_{1/2}(s)$ and
\eqref{lj:eq:Stieltjesone}. For the second use
$Z_2(s)=\eta(s-1)$,
$\eta(s)=(1-2^{1-s})\zeta(s)$,
$Z_2(0)=1/4$, and
$Z_2'(0)=-\log2/3-3\zeta'(-1)$ in
\eqref{lj:eq:Ijets}. The second order derivative is
\begin{equation}
 \int_{\mathbb R}\frac{\pi}{\cosh(x/2)}
       \left.\partial_s^2P_s(x)\right|_{s=1}\,\dd x
 =-2\pi\gamma_2'(1/2)-4\pi\log2\,\gamma_1'(1/2)
       +\pi^3\log^22.
 \label{lj:eq:exampleSecond}
\end{equation}
These formulas state the normalization explicitly; in particular, the
$2^{1-s}$ factor cannot be held constant during order differentiation.

\subsection{What this resolves and what remains}

Theorem~\ref{lj:thm:oddjets} resolves Q8 in the source's requested
Gamma--Stieltjes--$Z_\rho$-jet form. Theorem~\ref{lj:thm:Euler} supplies
a convergent evaluator and explicit harmonic coefficient calculus for
every reflected coordinate, while Theorem~\ref{lj:thm:polylog} transports
the result to ordinary Gamma-weighted polylogarithm integrals. This does
not assert that arbitrary noninteger-depth spectral jets reduce to a
finite algebra of classical zeta and Gamma values. Such a further
reduction would be a different arithmetic statement.

No false theorem was found in the inspected continuous-depth report.
Its caution that one cannot differentiate a formula known only at
integer $m$ is correct. The necessary addition is a holomorphic local
identity, \eqref{lj:eq:local}, together with the extra reflected jets.
The tilted report already treats independent shift and order derivatives
of its matched convolution family; those results are retained as prior
work and are not recounted as discoveries here.

Two concrete questions remain. First, determine whether selected rational
noninteger depths admit additional finite reductions of
$\partial_s^qZ_\rho(-2m)$, specifying the comparison algebra in advance.
Second, develop a cancellation-aware implementation of the Euler sums
with proved arithmetic error bounds. The analytic theorem supplies
normal convergence, while finite floating-point agreement by itself
does not provide a certificate for a transcendental identity.
